\section{第\ref{sec:dofaalt}章　\nt{DOPA}の他の理論}

拡張された描像はそれぞれ、ドーパミンの独自の直接測定（反転点のばらつき、不快なものと新しいものへの応答、ゆっくりした増加、味の変化への応答）に基づいている。しかしそれらを説明する式は理論であり、そのうち二つの間では実験がまだ互いに矛盾している。

{\footnotesize
\setlength{\tabcolsep}{3pt}\renewcommand{\arraystretch}{1.15}
\begin{longtable}{>{\raggedright}p{0.5cm}>{\raggedright}p{4.0cm}>{\raggedright}p{5.6cm}>{\raggedright}p{2.3cm}>{\raggedright\arraybackslash}p{1.7cm}}
\toprule
No. & 主張 & 実験と測定内容 & 出典 & 状態 \\
\midrule\endfirsthead
\toprule
No. & 主張 & 実験と測定内容 & 出典 & 状態 \\
\midrule\endhead
1 & 非対称な規則~\eqref{eq:asymmetric}は点~\eqref{eq:expectile}に収束する。$\kappa=1/2$ではそれは平均、確率$1/2$の1滴では$V=\kappa$（図~\ref{fig:expectiles}） & 点~\eqref{eq:expectile}はエクスペクタイル、つまり正と負の偏差に異なる重みをつけた最小二乗問題の解である。本章の例の導出は本章内にある。 & \cite{NeweyPowell1987}、本章で導出 & 理論 \\
2 & \nt{DOPA}ニューロンごとに反転点が異なり、非対称性の順に並ぶ（命題~\ref{prop:expectile}） & マウス。光遺伝学的に標識した腹側被蓋野の\nt{DOPA}ニューロン、大きさの異なる報酬。バーストが落ち込みに変わる点はニューロンごとに異なる。応答の非対称性が大きいニューロンほどその点は高い。群の応答から報酬分布の形が復元される。これがばらつきの主な機能であることは仮説。 & \cite{Dabney2020} & 測定済み \\
3 & \nt{DOPA}ニューロンの一部は不快なものにバーストで応答する & サル。ある\nt{DOPA}ニューロンは報酬の予告信号で興奮し、目への空気噴射の予告信号で抑制される。別のものは両方で興奮する。二つの群は中脳の異なる部分にある。 & \cite{Matsumoto2009, BrombergMartin2010} & 測定済み \\
4 & 誤差の絶対値または新奇性が学習率を決める & 引用した研究には、\nt{DOPA}ニューロンの重要性信号が学習率を変えることを示す直接の実験はない。レビューは「注意、重要」という信号の役割を提案している。 & \cite{BrombergMartin2010} & 仮説 \\
5 & 危険の軸のための別の配線 & マウス。線条体の後端に向かう\nt{DOPA}ニューロンは報酬ではなく新しい強い刺激に応答する。その破壊は新しい物体の回避を弱める。 & \cite{Menegas2018} & 測定済み \\
6 & 最適な行動時間$\tau_a^*=\sqrt{C/\bar R}$~\eqref{eq:vigor} & 和$C/\tau_a+\bar R\tau_a$の最小値。導出は本章。元のモデルでは、行動の速さは平均報酬に等しい時間のコストで決まる。 & \cite{Niv2007}、本章で導出 & 理論 \\
7 & ドーパミンのゆっくりしたレベルは$\bar R$を運び、行動の速さを決める（命題~\ref{prop:multiplex}） & ラット、側坐核でのマイクロダイアリシスとボルタンメトリー。ドーパミンのレベルは分単位で報酬の頻度と行動の速さに従う。ヒトではL-DOPAが反応時間の平均報酬頻度への依存を強める。遅いレベルがまさに$\bar R$に等しいことは証明されていない。 & \cite{Hamid2016, Beierholm2013vigor} & 仮説 \\
8 & ゆっくりしたレベルは二つの源からなる。放出はスパイクのレートを変えずに出力で変わるが、ゆっくりした増加はレート自体にもある & ラット、一つの課題。側坐核での放出は変化する報酬の期待に従うが、同定された腹側被蓋野\nt{DOPA}ニューロンのスパイクのレートは従わない。仮想通路のマウス（10行）。報酬への接近に伴うなめらかな増加は、同定された\nt{DOPA}ニューロンのスパイクにもある。 & \cite{Mohebi2019, Kim2020} & 測定済み \\
9 & 動物が遠くの報酬に近づく間、ドーパミンはなめらかに増える & 迷路のラット、線条体でのボルタンメトリー。ドーパミン濃度は目標に近づくにつれ数秒かけて上がり、その傾きは距離と報酬の大きさに依存する。 & \cite{Howe2013ramp, Hamid2016} & 測定済み \\
10 & 増加は推定の精密化に伴う誤差$\delta$である~\eqref{eq:ramp}。通路での跳躍はバーストを生む（図~\ref{fig:ramp}） & 式と毎秒$\delta=0.75\,V$という数値は本章で導出。実験：仮想通路のマウスを目標近くへ瞬時に移す。細胞体のスパイク、細胞体と出力のカルシウム、線条体での濃度は、期待$V$ではなく誤差のように応答する。二つの説明のどちらがすべての出力で正しいかは議論中。 & \cite{Kim2020} & 測定済み \\
11 & 振り返り：ドーパミンは因果的結びつきの尺度~\eqref{eq:retro}を伝える & マウス、この理論と誤差~\eqref{eq:ctd}が異なる予測をする実験での側坐核のドーパミン放出。いくつかの実験で放出は前者に従った。 & \cite{Jeong2022} & 仮説 \\
12 & 学習中のドーパミン応答は報酬から予告信号へ徐々に移る & マウス、学習の進行に伴う腹側線条体での\nt{DOPA}ニューロン出力の信号。バーストは\eqref{eq:ctd}による学習が要求するとおり、報酬の時点から予告信号の時点へ徐々に移る。 & \cite{Amo2022} & 測定済み \\
13 & \nt{DOPA}ニューロンは同じ価値での報酬の味の変化に応答する & ラット、腹側被蓋野ニューロンの記録。ジュースを価値の等しい別のものに替えると、誤差~\eqref{eq:ctd}はゼロなのに応答が生じる。 & \cite{Takahashi2017} & 測定済み \\
14 & 人工的なバーストは二つの中立な刺激の結びつきを教える & ラット、報酬なしに二つの刺激をあらかじめ対にする実験。第一から第二への移行時に\nt{DOPA}ニューロンを光遺伝学的に活性化すると結びつきが生じ、抑制するとその学習が妨げられる。 & \cite{Sharpe2017} & 測定済み \\
15 & 特徴ごとの誤差のベクトル~\eqref{eq:vectortd} & 特徴の予測誤差の理論。ベクトル形式の直接の検証はない。 & \cite{Gardner2018} & 仮説 \\
16 & 一つの課題の中で\nt{DOPA}ニューロンごとに異なる量を運ぶ & 仮想通路のマウス、\nt{DOPA}ニューロンの2光子イメージング。報酬に関係するもの、速度と旋回に関係するもの、予告信号と選択の正確さに関係するものがある。 & \cite{Engelhard2019} & 測定済み \\
17 & 配線ではなく配線の束（命題~\ref{prop:bundle}） & 2--16行のまとめ。各分解は個別に測定されている。それらがすべて一つの信号の部分であるという統一的な描像は実験で示されていない。 & --- & 仮説 \\
\bottomrule
\end{longtable}}
